
Il
existe un monomorphisme canonique de @@M0@@ dans le foncteur @@M1@@ des données de recollement sur @@M2@@. Par définition @@M3@@ est égal au foncteur @@M4@@ où les @@M5@@ sont les projections, qui est représenté par un schéma affine @@M6@@ de présentation finie sur @@M7@@ d'après 2.4. Il reste donc à représenter le monomorphisme
@@M8@@
par un morphisme de présentation finie, qui sera nécessairement quasi-affine (d'après EGA IV 8.11.2).

Remarquons en passant que si l'on y tient on peut restreindre @@M9@@ en imposant la condition de cocycle, mais ce n'est pas indispensable pour la suite.

Après le changement de base @@M10@@, on est donc ramené au problème suivant : si @@M11@@ est un faisceau inversible sur @@M12@@ muni d'une donnée de recollement @@M13@@, représenter le sous-foncteur de @@M14@@ «~qui rend cette donnée effective~». Mais on a le

\textbf{Lemme 4.3.} \uline{Sous les conditions de} 4.1\uline{, soit} @@M15@@ \uline{un faisceau inversible sur} @@M16@@ \uline{muni d'une donnée de recollement}
@@M17@@
\uline{Pour que cette donnée soit effective, il faut et il suffit que les cinq conditions suivantes soient réalisées.}
\begin{enumerate}[label=\arabic*)]
\item \uline{Pour tout} @@M18@@, @@M19@@ \uline{est plat sur} @@M20@@.
\item \uline{Pour tout} @@M21@@, @@M22@@ \uline{est plat sur} @@M23@@.
\item \uline{Le conoyau de} @@M24@@ \uline{est plat sur} @@M25@@\uline{, (où le dernier isomorphisme est} @@M26@@).
\item \uline{Le noyau} @@M27@@ \uline{de} @@M28@@ \uline{est inversible sur} @@M29@@.
\item \uline{Le morphisme composé} @@M30@@ \uline{est un isomorphisme.}
\end{enumerate}
